Both Marxer et al.~\cite{Marxer2023} and Luo et al.~\cite{Luo2023}
implement entangling gates by applying a baseband flux pulse to a
tunable transmon coupler. The pulse changes the interaction between
two data transmons and thereby controls their relative phases.
The relevant quantity is the \emph{conditional phase}: the phase of
a joint state after subtracting the phases attributable to the
individual systems. If a diagonal operation maps
$\ket{ij}\rightarrow e^{i\theta_{ij}}\ket{ij}$, the conditional
phase is
\begin{equation}
    \phi_{ij}
    = \theta_{ij}-\theta_{i0}-\theta_{0j}+\theta_{00}.
    \label{eq:conditional-phase}
\end{equation}
This quantity is invariant under local diagonal corrections and
therefore characterizes the entangling part of the operation.
A two-qubit CZ gate requires $\phi_{11}=\pi$, whereas the
generalized two-qutrit CZ gate requires a prescribed pattern of
four conditional phases.

\paragraph{Two-qubit CZ through a coherent excursion.}

In Marxer's implementation, the qubits are biased near the
$\ket{11}\leftrightarrow\ket{20}$ resonance. The coupler index is
omitted in this notation. Using the approximate transmon spectrum,
\begin{equation}
    \frac{E_{11}}{\hbar}\simeq\omega_1+\omega_2,
    \qquad
    \frac{E_{20}}{\hbar}\simeq2\omega_1+\alpha_1,
\end{equation}
where $\omega_1$ and $\omega_2$ are the qubit transition angular
frequencies and $\alpha_1<0$ is the first transmon's anharmonicity.
The corresponding detuning is
\begin{equation}
    \delta
    = \frac{E_{20}-E_{11}}{\hbar}
    \simeq\omega_1-\omega_2+\alpha_1.
\end{equation}
A flux pulse applied to the coupler activates the effective
exchange interaction while the qubit biases remain fixed.
After subtracting the energy of $\ket{11}$, a minimal model for
the relevant subspace is
\begin{equation}
    \frac{H_{\mathrm{sub}}(t)}{\hbar}
    =
    \begin{pmatrix}
        0 & G(t)\\
        G(t) & \delta(t)
    \end{pmatrix}_{\{\ket{11},\ket{20}\}},
    \qquad
    G\simeq\sqrt{2}\,g_{\mathrm{eff}}.
    \label{eq:two-state-model}
\end{equation}
The factor $\sqrt{2}$ arises from the transmon ladder-operator
matrix element. Coupler dressing can modify both the effective
coupling and detuning during the pulse.

For constant coupling at exact resonance, $\delta=0$, the
evolution from $\ket{11}$ is
\begin{equation}
    \ket{\psi(t)}
    = \cos(Gt)\ket{11}-i\sin(Gt)\ket{20}.
\end{equation}
At $T=\pi/G$, the population completes a round trip and
\begin{equation}
    \ket{\psi(T)}=-\ket{11}.
\end{equation}
Although the population returns to its initial value, the
amplitude acquires a phase of $\pi$. This distinction between
population return and amplitude return provides the essential
intuition for the gate: the auxiliary state $\ket{20}$ allows a
cyclic excursion that changes the phase of the computational
state $\ket{11}$.

For a constant detuned pulse, the same model gives
\begin{equation}
    \Omega=\sqrt{\delta^2+4G^2},
    \qquad
    T=\frac{2\pi}{\Omega},
    \qquad
    \ket{11}\rightarrow
    e^{i(\pi-\delta T/2)}\ket{11},
    \label{eq:detuned-return}
\end{equation}
in the chosen frame. Thus, population return alone does not
guarantee the desired conditional phase. The experimental pulse
is shaped and slightly detuned, so
Eqs.~\eqref{eq:two-state-model}--\eqref{eq:detuned-return}
provide intuition rather than a complete model of the waveform.
The pulse must be calibrated to suppress final leakage and
satisfy
\begin{equation}
    \theta_{11}-\theta_{10}-\theta_{01}+\theta_{00}
    =\pi \pmod{2\pi}.
\end{equation}
After correcting the local phases, the resulting operation is
\begin{equation}
    U_{\mathrm{CZ}}=\operatorname{diag}(1,1,1,-1).
\end{equation}
Marxer et al.\ implement this evolution using a Slepian-shaped
coupler pulse, with a $22\,\mathrm{ns}$ interaction time and a
$33\,\mathrm{ns}$ duration including padding~\cite{Marxer2023}.

\paragraph{Two-qutrit controlled phase through energy shifts.}

Luo's experiment uses $\ket{0}$, $\ket{1}$, and $\ket{2}$ of each
data transmon as logical states. An effective description follows
the dressed computational states as the coupler frequency changes.
Let $E_{ij}$ denote their energies, with the ground-state coupler
index suppressed, and define
\begin{equation}
    \chi_{ij}(t)
    =
    \frac{
        E_{ij}(t)-E_{i0}(t)-E_{0j}(t)+E_{00}(t)
    }{\hbar},
    \qquad i,j\in\{1,2\}.
    \label{eq:cross-kerr}
\end{equation}
After removing local energy contributions, the effective
conditional Hamiltonian is
\begin{equation}
    \frac{H_{\mathrm{cond}}(t)}{\hbar}
    =
    \sum_{i,j=1}^{2}
    \chi_{ij}(t)\proj{ij}.
\end{equation}
Within this diagonal description, the accumulated phases are
\begin{equation}
    \phi_{ij}
    =-\int_0^T\chi_{ij}(t)\,dt,
    \label{eq:integrated-phase}
\end{equation}
where the sign follows the convention
$U=\mathcal{T}\exp[-(i/\hbar)\int_0^T H(t)\,dt]$.

The physical origin of these conditional shifts is multilevel
hybridization: a joint state's energy differs from the sum of
the corresponding individual energies. Changing the coupler
frequency controls this nonadditivity. Pulse amplitude therefore
controls the phase rates, while pulse duration controls the
accumulated angles. Equation~\eqref{eq:integrated-phase} includes
the contributions from both the ramps and the plateau.
Experimentally, the interaction shifts are calibrated using
conditional Ramsey measurements, which compare phase evolution
for different states of the other qutrit~\cite{Luo2023}.

The target operation is the generalized qutrit CZ,
\begin{equation}
    U_{\mathrm{CZ}_3}\ket{ij}
    =e^{i2\pi ij/3}\ket{ij},
    \qquad i,j\in\{0,1,2\}.
\end{equation}
Its four nontrivial conditional phases satisfy
\begin{equation}
    \begin{pmatrix}
        \phi_{11} & \phi_{12}\\
        \phi_{21} & \phi_{22}
    \end{pmatrix}
    =
    \begin{pmatrix}
        2\pi/3 & 4\pi/3\\
        4\pi/3 & 2\pi/3
    \end{pmatrix}
    \pmod{2\pi}.
    \label{eq:qutrit-target}
\end{equation}

Luo et al.\ construct this pattern using two identical coupler
flux pulses, separated and followed by simultaneous
$\ket{1}\leftrightarrow\ket{2}$ microwave swaps on both qutrits.
The time-ordered sequence is
\begin{equation}
    \text{flux pulse}
    \;\longrightarrow\;
    \text{swaps}
    \;\longrightarrow\;
    \text{flux pulse}
    \;\longrightarrow\;
    \text{swaps}.
\end{equation}
Suppose one flux pulse contributes conditional phases
$(a,b,c,d)$ to
$(\ket{11},\ket{12},\ket{21},\ket{22})$.
For ideal swaps, and up to local phase corrections, the complete
sequence contributes
\begin{equation}
    (\phi_{11},\phi_{12},\phi_{21},\phi_{22})
    =(a+d,\;b+c,\;b+c,\;a+d).
    \label{eq:phase-pairing}
\end{equation}
For example, an input $\ket{11}$ samples the $\ket{11}$
interaction shift during the first pulse and the $\ket{22}$
shift during the second. An input $\ket{22}$ samples the same
shifts in reverse order. The swaps therefore enforce equality
within each pair and reduce the calibration to two phase sums:
\begin{equation}
    a+d=\frac{2\pi}{3},
    \qquad
    b+c=\frac{4\pi}{3}
    \pmod{2\pi}.
\end{equation}

In the paper's qutrit--coupler--qutrit notation, the resulting
operation on the dressed computational manifold is
\begin{equation}
\begin{aligned}
    U_{\mathrm{Cphase}}
    ={}&
    e^{i2\pi/3}
    \left(\proj{101}+\proj{202}\right)\\
    &+
    e^{-i2\pi/3}
    \left(\proj{102}+\proj{201}\right)\\
    &+
    \sum_{\substack{
        i,j\in\{0,1,2\}\\
        i=0\ \mathrm{or}\ j=0
    }}
    \proj{i0j}.
\end{aligned}
\end{equation}
The middle zero denotes the ground-state coupler label.
The final sum specifies that states with either qutrit in
$\ket{0}$ acquire no conditional phase after local corrections.
The equality $e^{-i2\pi/3}=e^{i4\pi/3}$ makes the connection to
Eq.~\eqref{eq:qutrit-target} explicit~\cite{Luo2023}.
